\section{ATOMIC y COMET: un Language Model no es un Knowledge Model}

Esta sección pasa de la evaluación a la representación. Un modelo de lenguaje general optimiza la continuación de texto, mientras que un commonsense system a menudo necesita responder typed questions: cuál es la causa de un evento, qué quieren los participantes, qué ocurrirá después. ATOMIC define el espacio de preguntas con symbolic relations, y COMET amplía las respuestas mediante neural generation.

\subsection{División del trabajo entre el Symbolic Graph y el Neural Model}

La slide coloca ATOMIC y COMET lado a lado: el primero es un symbolic commonsense knowledge graph escrito por personas mediante crowdsourcing; el segundo se entrena sobre esos relation-conditioned triples y se convierte en un neural commonsense model. No compiten entre sí: el graph aporta el schema y la supervisión, y el model aporta la generalización y la generación.

\lecturefigure{slide-13-language-models-not-knowledge-models.jpg}{Un language model y un knowledge model tienen objetivos distintos}{Diapositiva V013 recuperada de la grabación oficial; Stanford CS25 Lecture 16}

La página siguiente incorpora COMET-ATOMIC 2020 y destaca el ciclo entre el machine-readable graph y el generative model: primero las personas definen la relation ontology y el seed knowledge; después el modelo aprende $p(t\mid h,r)$ y genera tail inferences para eventos nuevos. Aquí $h$ es el head event, $r$ es el tipo de relación y $t$ es la commonsense conclusion generada.

\lecturefigure{slide-14-atomic-comet-human-symbolic-neural.jpg}{La relación symbolic-neural entre ATOMIC, COMET y COMET-ATOMIC 2020}{Diapositiva V014 recuperada de la grabación oficial; Hwang et al., 2021}

Formalmente, el objetivo de entrenamiento de COMET puede escribirse como
\begin{equation}
\mathcal{L}_{\mathrm{COMET}}(\theta)
=-\mathbb{E}_{(h,r,t)\sim\mathcal{D}}
\log p_{\theta}(t\mid h,r).
\end{equation}
$\mathcal{D}$ son los datos de tripletas de conocimiento; $h$ y $r$ delimitan de manera explícita qué tipo de inferencia debe generarse. Frente a una unrestricted continuation, esta interface se parece más a «consultar un modelo de conocimiento».

\subsection{Los Relation Types son el sistema de coordenadas del conocimiento}

El grafo completo de ATOMIC despliega relation types como social interaction, physical entity y event-centered. Al leer la figura no conviene memorizar nodo por nodo, sino observar la semántica de los colores y de las aristas: un mismo evento puede conectarse con agent intent, agent reaction, other reaction, precondition y effect. El relation type convierte las «inferencias posibles» en una tarea descomponible y permite que el evaluator revise cada capacidad por separado.

\lecturefigure{slide-15-atomic2020-full-graph.jpg}{ATOMIC 2020: grafo de conocimiento con nodos de eventos y 23 tipos de relaciones de sentido común}{Diapositiva V015 recuperada de la grabación oficial; Hwang et al., 2021}

\begin{knowledgebox}{Por qué importa la Typed Relation}
Si solo se deja que el modelo continúe el texto libremente, lo «razonable» puede ir en un sinnúmero de direcciones. Al especificar \texttt{xIntent}, \texttt{xEffect}, \texttt{oReact} o una physical relation, la recolección de datos, el entrenamiento y el análisis de errores obtienen coordenadas. El costo es que la propia ontology determina lo que el sistema puede expresar: las relaciones que no se modelan no aparecen por sí solas.
\end{knowledgebox}

\subsection{Los límites de la comparación entre COMET, GPT-3 y GPT-2}

La slide comparativa presenta la human-judged commonsense quality: COMET-BART, GPT-3 few-shot y GPT-2 zero-shot tienen configuraciones distintas. Lo primero que debe leerse son las etiquetas de método, no qué número es mayor; COMET utiliza datos y supervisión específicos, GPT-3 depende de few-shot prompting y GPT-2 es más pequeño y zero-shot. El resultado respalda que «un modelo de conocimiento especializado puede ser más pequeño y más útil», pero no permite concluir una relación causal atribuible solo a la architecture.

\lecturefigure{slide-16-comet-vs-gpt3-gpt2.jpg}{Comparación human-judged commonsense entre COMET y los modelos de lenguaje generales}{Diapositiva V016 recuperada de la grabación oficial; Stanford CS25 Lecture 16}

\teachervoice{La ponente aclara por iniciativa propia que no se trata de una apple-to-apple comparison. Reconoce que el scale jump de GPT-2 a GPT-3 es interesante, pero también señala que GPT-3 resulta demasiado grande para muchos ingenieros e investigadores. El énfasis de la clase está en la system utility: un modelo más pequeño, consultable y que pueda integrarse en sistemas posteriores puede tener más valor.}

\subsection{El valor de reutilizar un Knowledge Backbone}

La última página coloca COMET / ATOMIC en el centro y lo conecta con persona-aware conversation, figurative language, storytelling, fantasy games e interactive learning. Lo que muestra no es que todas las aplicaciones estén «resueltas», sino que una relation interface unificada reduce el costo de construir una y otra vez componentes de sentido común. El desempeño de las aplicaciones sigue limitado por la coverage, la quality y el cultural scope.

\lecturefigure{slide-17-commonsense-backbone-applications.jpg}{ATOMIC / COMET como commonsense backbone de varias tareas posteriores}{Diapositiva V017 recuperada de la grabación oficial; Stanford CS25 Lecture 16}

\subsection{Resumen de la sección}

ATOMIC formula los problemas de sentido común como typed symbolic queries, y COMET condensa la supervisión del grafo en un generative model. La diferencia entre un modelo de lenguaje y un modelo de conocimiento no está en si usa Transformer, sino en el objetivo de entrenamiento, la interfaz y la forma de validación. La siguiente sección plantea otra pregunta: si escribir grafos de conocimiento a mano es demasiado lento, ¿puede un modelo grande producir el corpus y luego usarse un critic para destilarlo en un modelo más pequeño?
